\documentclass[11pt,a4paper]{ctexart}
\usepackage[margin=23mm]{geometry}
\usepackage{amsmath,amssymb,fancyhdr}
\usepackage[colorlinks=true,linkcolor=blue,pdftitle={hw 10-25}]{hyperref}
\pagestyle{fancy}\fancyhf{}
\fancyhead[L]{微分方程}\fancyhead[R]{hw 10-25}
\fancyfoot[C]{\thepage}\setlength{\headheight}{15pt}
\setlength{\parindent}{0pt}\setlength{\parskip}{5pt}
\newcommand{\xx}{\boldsymbol{x}}
\newcommand{\bb}{\boldsymbol{b}}
\begin{document}
\begin{center}{\LARGE\bfseries hw 10-25}\\[7pt]题目与解答\end{center}
\textbf{范围：}第 93 页第 9 题 (1)；第 94 页第 10 题。题干依据用户提供的截图，页码沿用用户标注。

\section*{第 93 页，第 9 题 (1)}
\textbf{题目。}求解下列非齐次线性系统的初值问题：
\[
\text{(1)}\quad
\begin{cases}
\displaystyle\frac{dx_1}{dt}=-\frac2t x_1+1,\quad
\displaystyle\frac{dx_2}{dt}=\left(1+\frac2t\right)x_1+x_2-1
\quad(t>0),\\[7pt]
\displaystyle x_1(1)=\frac13,\quad x_2(1)=-\frac13.
\end{cases}
\]

\textbf{解答。}先解第一个方程：
\[
x_1'+\frac2t x_1=1.
\]
在 $t>0$ 上，积分因子为 $t^2$，故
\[
(t^2x_1)'=t^2.
\]
从 $1$ 到 $t$ 积分，利用初值 $x_1(1)=1/3$，得到
\[
t^2x_1(t)-\frac13=\frac{t^3-1}{3},
\qquad\text{即}\qquad x_1(t)=\frac t3.
\]

将原系统的两个方程相加，有
\[
(x_1+x_2)'=
-\frac2t x_1+1+\left(1+\frac2t\right)x_1+x_2-1
=x_1+x_2.
\]
令 $y=x_1+x_2$，则
\[
y'=y,\qquad y(1)=\frac13-\frac13=0.
\]
因此 $y(t)=y(1)e^{t-1}=0$，从而 $x_2(t)=-x_1(t)=-t/3$。
所求解为
\[
\boxed{\quad x_1(t)=\frac t3,\qquad x_2(t)=-\frac t3,\qquad t>0.\quad}
\]
系统的系数和非齐次项在 $(0,\infty)$ 上连续，由线性系统初值问题的存在唯一性定理，此解是该区间上满足所给初值的唯一解。

\newpage
\section*{第 94 页，第 10 题}
\textbf{题目。}设 $A(t)$ 和 $b(t)$ 分别是区间 $I$ 上的连续 $n$ 阶矩阵函数和 $n$ 维向量函数，并且 $b(t)\ne0$，证明：非齐次线性系统
\[
\frac{dx}{dt}=A(t)x+b(t)
\]
在区间 $I$ 上至多有 $n+1$ 个线性无关的解。

\textbf{证明。}这里的线性无关是指作为区间 $I$ 上的向量函数线性无关，即线性组合中的系数为与 $t$ 无关的常数。

任取上述非齐次系统的 $n+2$ 个解
\[
\xx^{(1)}(t),\ \xx^{(2)}(t),\ \ldots,\ \xx^{(n+2)}(t).
\]
只需证明它们线性相关。固定 $t_0\in I$，在 $\mathbb R^{n+1}$（复系数情形则在 $\mathbb C^{n+1}$）中考虑 $n+2$ 个向量
\[
\begin{pmatrix}\xx^{(1)}(t_0)\\1\end{pmatrix},\quad
\begin{pmatrix}\xx^{(2)}(t_0)\\1\end{pmatrix},\quad\ldots,\quad
\begin{pmatrix}\xx^{(n+2)}(t_0)\\1\end{pmatrix}.
\]
它们必线性相关。因此存在不全为零的常数 $c_1,\ldots,c_{n+2}$，使
\[
\sum_{j=1}^{n+2}c_j\xx^{(j)}(t_0)=0,
\qquad
\sum_{j=1}^{n+2}c_j=0.
\]
令
\[
\boldsymbol{y}(t)=\sum_{j=1}^{n+2}c_j\xx^{(j)}(t).
\]
由于每个 $\xx^{(j)}$ 都满足原系统，得到
\begin{align*}
\boldsymbol{y}'(t)
&=\sum_{j=1}^{n+2}c_j\bigl(A(t)\xx^{(j)}(t)+\bb(t)\bigr)\\
&=A(t)\boldsymbol{y}(t)+\left(\sum_{j=1}^{n+2}c_j\right)\bb(t)
=A(t)\boldsymbol{y}(t).
\end{align*}
同时 $\boldsymbol{y}(t_0)=0$。

由齐次线性系统初值问题的唯一性，$\boldsymbol{y}$ 与零解有相同初值，故
\[
\boldsymbol{y}(t)\equiv0\qquad(t\in I).
\]
于是
\[
\sum_{j=1}^{n+2}c_j\xx^{(j)}(t)\equiv0,
\]
且这些常数不全为零，说明任意 $n+2$ 个解都线性相关。因此
\[
\boxed{\text{该非齐次线性系统在 }I\text{ 上至多有 }n+1\text{ 个线性无关的解。}}
\]
\end{document}
